% chapters/ch06.tex —— 第 6 章
% 体例：具体失败 → 设计问题 → 机制解释 → 对照实验 → 条件变化 → 选择原则
% 数字身份纪律：本章外部数字全部来自方法卡所登记的原文位置，一律 \cited{}；
% 两个由引用数字算出的量（延迟降幅、增益相等）用 \calcnum{}。本章无新的模型实测；离线接口检查与教学算例分别标识。
\chapter{记忆操作的学习化：从规则到策略}
\label{ch06}

\chaptergoal{谁来写、留不留、取不取——这三件事能不能学？}
\chapterprops{P3关注形式与容量如何共同影响迁移，本章从操作机制展开这一问题，给出学习对象的可检查判据与统一对照表，供E03b等后续设计使用。}
\chapterfig{图~\ref{fig:ch06-ladder}：操作选择、环境校验与延迟反馈的边界}
\chapterpitfall{动作空间里有“删除”，评测维度里却没有冲突消解}

% ---------- 1. 具体失败 ----------
\begin{failurebox}
一批长期记忆系统里，有一类的召回率恰好是 \cited{1.000}。
在该论文整理的 LoCoMo 候选记忆测试集（\cited{N=225}）上，A-mem 的精确率 \cited{0.371}、召回 \cited{1.000}、
F1 \cited{0.541}、单次延迟 \cited{3831} 毫秒\cite{amac-2603}。

这里评估的是候选记忆的准入分类，精确率与召回率相对于标注的准入标签定义，
不是后续任务真正采用经验的比例。召回为一表示标注应准入的候选没有被漏掉，
并不单独证明所有候选都被保留，更不能由这个数反推出系统的全部写入机制。
低精确率提示过度准入，需要进一步检查额外条目对检索和任务结果的影响。


同一批测试集上，A-MAC 把准入显式做成一个打分与阈值的决策，
F1 提到 \cited{0.583}、延迟降到 \cited{2644} 毫秒（相对 A-mem 降低 \calcnum{31.0\%}）\cite{amac-2603}。
分数的差别有意思，但本章关心的不是谁高：
\textbf{两个系统里“留不留”这个决策的做出者完全不同}——
A-Mem 的主要机制是笔记与关系维护，A-MAC 则显式学习五个特征的权重与准入阈值。
这两个工作关注的环节不同，准入表中的差异不能直接当成完整记忆系统的优劣排名。
\end{failurebox}

\designquestion{“记忆操作的学习化”判据不是“有没有用 RL”，而是：写、留、取、删这四个决策里，
有几个的决策权从设计者手里交给了学习过程。}

第 5 章讨论了“经验写成什么样才可能跨任务”，本章追问它的前半段：
\textbf{这些操作本身由什么决定}。
固定规则与学习策略都受观测、动作接口和评价目标约束。
学习可以适应规则不易逐一覆盖的模式，也可能利用奖励缺口；规则也可以稳定地满足某些已知约束。本章把两头的账都摆出来。

% ---------- 2. 机制解释 ----------
\section{准入：把“留不留”写成可检查的决策}
\label{sec:ch06-rule}

A-MAC 的做法是把准入从生成的副产品变成一次显式判决。
每个候选记忆 $m$ 先被切成语义上不可再分的单元，再做时间表达式与指代的消解，
最后算一个标量分\cite{amac-2603}：

\begin{equation}
S(m)=w_1 U(m)+w_2 C(m)+w_3 N(m)+w_4 R(m)+w_5 T(m),\qquad w_i\ge 0,\ \sum_{i=1}^{5}w_i=1,
\label{eq:ch06-amac}
\end{equation}

其中 $U$ 是未来效用、$C$ 是事实可信度、$N$ 是语义新颖性、$R$ 是时间鲜活性、$T$ 是内容类型先验。
在这五项打分中，$U$使用一次语言模型调用，温度设为0并缓存重复候选（温度零不保证外部服务逐字可复现）；
其余四项全部用可审计的规则算\cite{amac-2603}。候选在 $S(m)\ge\theta$ 时被准入（报告取 $\theta=\cited{0.55}$）。

\begin{table}[htbp]
\centering
\caption{论文整理的 LoCoMo 候选记忆测试集：准入精确率、召回率与延迟}
\label{tab:ch06-locomo}
\begin{tabular}{>{\raggedright\arraybackslash}p{3.0cm}cccc}
\toprule
方法 & 精确率 & 召回 & F1 & 延迟（毫秒） \\
\midrule
Random & \cited{0.278} & \cited{0.278} & \cited{0.278} & \cited{$<$1} \\
MemGPT & \cited{0.316} & \cited{0.333} & \cited{0.324} & \cited{2765} \\
MemoryBank & \cited{0.368} & \cited{0.583} & \cited{0.452} & \cited{2843} \\
等权五因子 & \cited{0.362} & \cited{0.694} & \cited{0.476} & \cited{2916} \\
A-mem & \cited{0.371} & \cited{1.000} & \cited{0.541} & \cited{3831} \\
A-MAC & \cited{0.417} & \cited{0.972} & \cited{0.583} & \cited{2644} \\
\bottomrule
\end{tabular}

\begin{minipage}{0.92\linewidth}
\vspace{2mm}
{\footnotesize 表注：全部数字取自 \cite{amac-2603} Table 1，条件为 LoCoMo 测试集 $N=\cited{225}$；
延迟含语言模型接口调用。等权五因子指式~\eqref{eq:ch06-amac} 取 $w_i=1/5$ 的对照。}
\end{minipage}
\end{table}

这张表里藏着一条容易被读漏的观察。
把 $T$（内容类型先验）这一项去掉，F1 从 \cited{0.583} 掉到 \cited{0.476}，
$\Delta F1=\cited{-0.107}$，是五项里最大的一条；
而等权基线恰好也是 \cited{0.476}\cite{amac-2603}。
也就是说，\textbf{在这份报告的数字上，“学出权重”带来的全部增益
（\calcnum{0.583-0.476=0.107}）与“移除内容类型先验的消融差值”在数值上相等}。

这不构成因果分解——两次消融走的是不同的对照路径。
这两条差值相等，不能证明权重优化的收益全部由类型先验贡献，更不能由消融降幅推断该项权重最大。
可以说，在原报告的消融设置中，移除类型先验带来的 F1 降幅最大。
权重与阈值是在标注数据上通过五折交叉验证和网格搜索选择的，不是强化学习。
换领域后应重新检查特征和阈值是否适用，不能预先断言先验一定失效或一定无需调整。


\subsection{从五项分数回到实际可检查的操作}

先辨清表中的样本是什么。原论文从\cited{30}段对话整理出约\cited{1{,}500}条带准入标签的候选记忆，报告训练、验证和测试比例为\cited{70\%}/\cited{15\%}/\cited{15\%}，表 1 的\cited{225}指测试候选数量\cite{amac-2603}。它不是整套原始问答基准的题数，也不是这些记忆已经改善下游回答的次数。公开参考代码的数据加载器先按对话划分，再展开候选\cite{amac-code-2026}；测试集内部的候选仍可能相关，不能仅凭候选条数就假定独立。原表的具体预测仍需用对应候选清单与运行记录核对。

五个名称还需要展开成可操作定义。原文的事实可信度使用候选与支持话语之间的最大 ROUGE-L 相似度\cite{amac-2603}。这个量衡量文字对齐，不能直接当作真实概率。一个几乎照抄原句但改变否定词的候选，可以保留大量重合；一句正确的改写也可能只保留较少字面重合。本书提出这些反例，是为了检查代理指标的边界，不是在没有实际运行的情况下声称某个具体候选得了多少分。将分数命名为“可信度”，不会额外赋予它事实验证能力。

新颖性需要同时读公式与实现。原文定义为一减去最大余弦相似度，并将五项特征描述为零到一之间\cite{amac-2603}。若余弦允许负值，未截断的公式本身不保证这个范围：本书构造向量$(1,0)$与$(-1,0)$，余弦为\calcnum{-1}，原始差值为\calcnum{2}。

公开参考代码补充了边界约定\cite{amac-code-2026}：空库返回\cited{1.0}；非空时在余弦可计算的前提下，将差值截断到\cited{[0,1]}。上述向量例子在该实现中返回\calcnum{1}，这里的原始差值与实际返回值必须区分。截断满足了范围约定，也使不同的负最大相似度落到同一上界；复现时应保留这项转换，而不是只抄论文中的未截断公式。

最后检查准入之后的冲突路径。原文叙述强调保留分数较高的表述，但算法 1 的分支写法更具体：当存在冲突且新候选分数更高时合并，否则加入新候选\cite{amac-2603}。按该伪代码直读，已通过阈值但不高于旧条目的冲突候选，会落入加入分支。公开参考代码的 \texttt{scorer.py} 在未注入特征提取器时采用占位分数。其 \texttt{admit} 接口返回准入决定、分数和特征，未在该接口中执行冲突合并\cite{amac-code-2026}，因此仍需由实际写入路径与运行轨迹判断如何处置这类候选。准入评分与冲突写入是不同接口，前者的代码不能证明后者采用了哪条分支。


\paragraph{不同表能否合并，需要先核对指标口径。}
同一报告的表 5 将测试候选分成 Personal 的\cited{127}条和 Professional 的\cited{98}条，F1分别为\cited{0.482}与\cited{0.338}，其等权均值列为\cited{0.410}；表 1 的总体 F1 则是\cited{0.583}\cite{amac-2603}。两表不能在缺少进一步说明时直接拼成一次完整的分域分析。

原因可以由计数检查。若各域使用同一套二分类预测与真值，令$D_g=2\mathrm{TP}_g+\mathrm{FP}_g+\mathrm{FN}_g$，其中三项分别为真正例、假正例与假负例数，且$D_g>0$。合并计数后的F1满足
\[
 F_{1,\mathrm{pooled}}=\frac{\sum_gD_gF_{1,g}}{\sum_gD_g}.
\]
它是按这些分母加权的平均，必在分域F1的最小值与最大值之间；权重也不是简单的样本条数。因而在上述共同口径假设下，\cited{0.583}不能由\cited{0.482}与\cited{0.338}合并得到，三位小数的舍入也不足以解释差距。如果两表采用不同预测、不同平均方式或不同运行，须先补齐这些信息。当前未拿到能核对这一步的原始分类计数，本书保留原表数字及该缺口，不替作者选择修正值，也不以此推断整项工作的结果均无效。

若把表间差距归因于两次不同的准入判定，至少要有多少候选改变判定？设两表使用同一批候选、逐条使用相同的二元标签，在总体和各域内分别按正类计数计算指标，再四舍五入到三位小数。满足表~\ref{tab:ch06-locomo}中 A-MAC 的精确率、召回率和F1的整数计数只允许\calcnum{36}或\calcnum{72}个正例。在两域内枚举与原论文表~5相容的计数，两种情形下改变判定的候选分别至少有\calcnum{36}个和\calcnum{24}个。例如，在后者中，总体表可以有\calcnum{70}个真正例和\calcnum{98}个假正例，分域表合起来可以有\calcnum{47}个真正例和\calcnum{99}个假正例。把两次判定尽量对齐，仍要漏掉原先准入的\calcnum{23}个正例，并额外纳入\calcnum{1}个负例。这个\calcnum{24}是给定口径下可以达到的下界。按对话先算后平均时需另行推导。要识别实际发生的判定变化与冲突合并，还需分别核对逐条预测和写入记录。

这些检查不会抹掉准入分类的价值。它们让读者知道，权重学习可以改变五项信号如何组合，却不能自动修正信号的语义、定义域和冲突处理。先把这些条件固定下来，再比较规则、权重优化与端到端策略，才能判断收益来自哪一层。

\section{组织：用模型维护笔记与关系}
\label{sec:ch06-selforg}

A-mem 走的是另一条路：不给准入打分，而是让记忆自己组织起来。
每条笔记是 $m_i=\{c_i,t_i,K_i,G_i,X_i,e_i,L_i\}$，
其中关键词、标签与上下文描述由语言模型按模板生成，向量 $e_i$ 由编码器算出；
新笔记写入时按余弦相似度取近邻，再由语言模型判断哪些笔记应建立链接、哪些近邻需要更新描述、关键词或标签\cite{amem-2502}。

这条路径的代价有两本账。
一本是 token 账：每次记忆操作约 \cited{1,200} 个 token，
对照基线报告约 \cited{16,900} 个 token，原文另汇总了 \cited{85--93\%} 的降幅范围\cite{amem-2502}。
这是成本段的概述；其表 1 的每题平均 token 长度还随模型配置变化，不能视作每一步的固定收费，具体口径见第 5 章。
同时，检索数量并非各模型各类别都为十：原文表 8 的 GPT-4o-mini 使用 \cited{40} 或 \cited{50} 条，
取决于问题类别。方法比较应保留这些实际配置及调参过程，不能用默认值代替报告成绩所用的设置。
另一本是规模账：候选检索、链接生成、近邻改写与索引维护随存储增长怎样变化。
D-MEM 将其对照方案描述为存在二次增长瓶颈，并报告在其带噪对话设置中降低接口token用量约\cited{80\%}\cite{dmem-2603}。
这是一项具体实现与负载下的报告，不能从“每次改写近邻”直接推出所有 A-Mem 实现的单次写入都是 $O(N^2)$。
应分别核对单次操作、累计过程、检索实现和固定近邻数；渐近复杂度也不等于实测延迟。 D-MEM公开归档的缓存条件与摄入费用覆盖，见第\ref{sec:ch13-gate}节。

把两节放在一起看，准入评分与关系维护的分界不在“用没用语言模型”：
A-MAC 也调用了一次语言模型给 $U$ 打分。
分界在于\textbf{决策的触发者与判决形式}——
A-MAC 是每个候选都要过一次阈值，A-mem 在写入时检查近邻，由模型逐项判断是否需要改写。
前者的准入分数可沿特征与阈值追查；后者需要从笔记、链接与演化请求追查。两种审计对象不同。

\section{触发：让记忆演化经过一道门}
\label{sec:ch06-gate}

D-MEM用固定路由规则控制是否进入深度演化，把演化触发变成了一道可观测的门。
一个轻量的评论家路由器结合语义新颖程度与长期效用，
将组合分数命名为智能体版奖励预测误差（RPE）；这里是生物启发的触发门分数，不是标准时序差分误差；
常规路由中，低于低阈值的输入跳过主流程，达到低阈值而未达到高阈值的输入只构建笔记并进入 $O(1)$ 的快速缓冲。达到高阈值的输入，例如事实矛盾或偏好反转，才触发慢速的 $O(N)$ 深度演化\cite{dmem-2603}。
跳过主流程的原始话语还会另存于 FIFO 影子缓冲，在常规记忆检索置信不足时作为回退上下文\cite{dmem-2603}。

这一支在本章的作用是对照。
如果后面那些用强化学习学出来的策略，学到的主要内容其实就是“少演化几次”，
就应加入相近触发率的规则门对照，检查学习策略是否提供额外收益，不能预先假定能省下多少训练。
它同时给出了一个可测的中间量：\textbf{演化触发率}。
任何声称学会了“何时更新”的系统，都应该能报出这个数，
否则无法区分“学会了挑时机”与“整体变懒了”。

\section{操作策略：在给定工具集上学习}
\label{sec:ch06-mdp}

直接训练操作选择的一类方法，做法是把记忆操作写进动作空间，
让策略自己决定调用哪一个。以 AgeMem 的形式化为例：
状态 $s_t=(C_t,M_t,T)$，动作 $a_t$ 同时包含语言生成与记忆操作，
目标是 $\theta^\ast=\arg\max_\theta \mathbb{E}_{\tau\sim\pi_\theta}[R(\tau)]$\cite{agemem-2601}。

\begin{table}[htbp]
\centering
\caption{不同方法的决策权归属：四个决策分别由谁做出}
\label{tab:ch06-who}
\begin{tabular}{>{\raggedright\arraybackslash}p{2.1cm}>{\raggedright\arraybackslash}p{2.3cm}>{\raggedright\arraybackslash}p{2.3cm}>{\raggedright\arraybackslash}p{2.3cm}>{\raggedright\arraybackslash}p{4.4cm}}
\toprule
方法 & 写 & 留 & 取 & 删 / 遗忘 \\
\midrule
A-MAC & 候选预处理流程 & \textbf{学习出的权重}过阈值 & 检索器 & 冲突合并规则 \\
A-mem & 写入事件 & 笔记构建与演化 & 余弦近邻 & 提示含prune；删除对象未明 \\
D-MEM & RPE Gate & 分级存储及缓冲留存 & 检索与缓冲回退 & 未列显式删除工具 \\
AgeMem & \textbf{RL 策略} & \textbf{RL 策略} & 策略发起，工具匹配 & DELETE 在工具集内 \\
Mem-$\alpha$ & \textbf{RL 策略} & \textbf{RL 策略} & 解耦 RAG 管线 & DELETE 在动作集内 \\
MEM1 & 策略生成 $<$\!IS$>$ & 上下文至多 2 个 $<$\!IS$>$ & 读取内部状态 & 框架清理旧轮，策略选择保留内容 \\
JitRL & 存三元组（固定规则） & 跨 episode 累积（固定规则） & 固定相似度检索 & 未定位显式遗忘 \\
\bottomrule
\end{tabular}

\begin{minipage}{0.92\linewidth}
\vspace{2mm}
{\footnotesize 表注：决策权归属依据各方法的原文机制描述\cite{amac-2603,amem-2502,dmem-2603,agemem-2601,memalpha-2509,mem1-2506,jitrl-2601}；
表中分别登记工具与模板已说明的动作，不据此保证记忆不会减少。A-Mem附录B.3的JSON示例含\texttt{prune}，正文说明主要描述连接强化与近邻更新；该模板没有明确删节点、删链接还是其他裁剪，需进一步核对实现。}
\end{minipage}
\end{table}

下面四个例子分别说明操作选择、写入策略、内部状态和潜记忆的学习对象：

\textbf{AgeMem} 把记忆管理做成六个工具——ADD、UPDATE、DELETE、RETRIEVE、SUMMARY、FILTER——
前三个工具改变长时记忆 $M_t$，后三个管理工作上下文 $C_t$。例如，策略决定何时调用 RETRIEVE 以及查询什么，工具按语义相似度从长时记忆取回条目，送入工作上下文。
训练分三阶段：先学长时记忆构建，再在干扰项下学上下文控制，最后联合推理；
由于记忆操作的收益要到任务结束才看得出，它把终端优势广播到同一轨迹的先前步骤\cite{agemem-2601}。
在五个长程基准上（只在 HotpotQA 上训练，其余零样本），
论文采用的跨数据集平均表现为 \cited{41.96\%}（Qwen2.5-7B-Instruct）与 \cited{54.31\%}（Qwen3-4B-Instruct），
相对去掉 RL 的自身版本分别高 \cited{8.53} 与 \cited{8.72} 个百分点\cite{agemem-2601}。
这组数字是本章最重要的一条证据：\textbf{它在相同工具框架下报告了训练版本与未训练版本的差异}，
工具集仍是人工固定的，训练还可能改变推理和其他行为，不能将全部增益只归因于调用时机。

\textbf{Mem-$\alpha$} 把写入动作展开成一个序列 $(a_1,\dots,a_{K_t})$，
动作集是 insert、update、delete，参数含记录号、记忆类型与内容。情景与语义记忆支持这三种操作，核心记忆只允许整体更新。
奖励由四项组成：下游问答正确率 $r_1$、工具格式 $r_2$、压缩率 $r_3$、内容合法性 $r_4$，
原文两处默认权重说明不一致，详见本章的数字核对说明；
策略采用 GRPO；原文将去掉 KL 项作为鼓励探索的设计，公开主训练脚本的配置差异见第~\ref{sec:ch06-config-identity} 节\cite{memalpha-2509,memalpha-code-2025}。
原文 Table 3 的七个任务列里，训练版本均高于同型号未微调基座\cite{memalpha-2509}。验证集均值经作者更正为\cited{0.619}，另一个测试任务集合的均值为\cited{0.592}\cite{memalpha-erratum-2025}。
原文把记忆架构（核心 / 情景 / 语义三组件）与 RL 解耦，
因此它学到的严格说是“\textbf{在固定结构上的写入策略}”，不是结构本身。

\textbf{MEM1} 取消了独立的记忆模块：每一步生成一段内部状态 $<$\!IS$_t>$ 来汇总历史并推理，
下一步把它与查询、环境反馈合并成新的 $<$\!IS$_{t+1}>$，
上下文里至多保留 2 个 $<$\!IS$>$、2 个查询与 1 条环境信息，
内部状态与行动生成由 PPO 训练\cite{mem1-2506}。
模型从上一步的内部状态接续推理，无需单设记忆检索器。遇到检索增强问答任务时，它仍可查询外部资料，查询结果作为环境反馈进入下一步。它学习的是任务内状态管理，不能仅凭窗口滚动就认定存在跨任务自进化。
代价写在训练侧：需要 \cited{4} 张 H100 或 H200 做 PPO，
报告采用训练后模型进行推理，不能把任务内状态变化等同于部署后的持续参数学习\cite{mem1-2506}。

\textbf{MemGen} 把“取”这件事也交了出去。
其形式化描述是在 token 生成过程中，根据当前隐藏状态计算调用概率
$p_j=\sigma(\mathrm{T_{trigger}}(h_{t,1},\dots,h_{t,j-1}))$，
并按伯努利采样决定 invoke 还是 skip。实际实现采用句级激活，只在指定分隔 token 处评估触发，其余位置的调用概率为零\cite{memgen-2509}。
一旦 invoke，编织器按同一状态生成一段定长 $K$ 的潜记忆
$M_t\in\mathbb{R}^{K\times d}$，把它拼回隐藏状态后再继续生成\cite{memgen-2509}。触发器与编织器分别由 LoRA 适配器实现，基础推理模型保持冻结，训练更新的是适配器参数\cite{memgen-2509}。
其原文报告跨多个基准的收益，比较条件各不相同；本章不以这些最大增益排列方法优劣\cite{memgen-2509}。
注意它的记忆不是自然语言条目，因此“这条记忆写了什么”无法直接审计。

\section{一支容易被放错位置的方法}
\label{sec:ch06-jitrl}

JitRL 说明了为什么需要区分任务策略学习与记忆操作学习。按本章的判据，它没有训练记忆操作策略。

它把经验写成三元组 $(s_i,a_i,G_i)$ 存进一个非参经验记忆 $M$，
检索时取 $k$ 近邻估计价值与优势，
再按带 KL 惩罚的策略优化闭式解把优势加到冻结基座的 logit 上：
$z'(s,a)=z(s,a)+\beta\,\hat A(s,a)$\cite{jitrl-2601}。
数学上的logit更新还需对应接口。原文第5.1节区分两种实现：读取候选编号token的对数概率，或将黑箱模型口头报告的置信分数变换为候选分数。后一种并未直接取得基础模型的完整logit向量；复现时应保留所用实现与候选集合。
这里被学习的是\textbf{任务策略的调制}，不是记忆操作：
写入是“存三元组”这条固定规则，检索是固定的相似度计算，
本书在原文中尚未定位到JitRL经验库的显式遗忘机制\cite{jitrl-2601}。该文附录E所述FIFO截断属于Memory基线，不能据此补成JitRL自身的操作。

JitRL利用经验调整任务策略，却没有学习存取操作的选择；A-MAC则确实根据标注数据拟合准入权重与阈值。
两者都可以不使用端到端的记忆操作强化学习，但学习对象不同，不能因此一并归为“操作全由人定”。
区分学习对象以后，换基座模型时需要核对的环节也就清楚了：
MEM1 的已训练状态管理行为依赖策略参数，换基座需要重新验证；
JitRL 的经验、状态表示、动作对应与优势尺度同样可能需要适配，不能保证只移动偏置就能迁移。

% figures/ch06-ladder.tex —— 第 6 章核心图（全书引用入口）
% 用法：在 book/chapters/ch06.tex 中 % figures/ch06-ladder.tex —— 第 6 章核心图（全书引用入口）
% 用法：在 book/chapters/ch06.tex 中 % figures/ch06-ladder.tex —— 第 6 章核心图（全书引用入口）
% 用法：在 book/chapters/ch06.tex 中 \input{figures/ch06-ladder}
% 图主体在 ch06-ladder-body.tex，便于单独编译与复用。
\begin{figure}[htbp]
\centering
\input{figures/ch06-ladder-body}
\caption{记忆操作的决策边界：策略提出动作，环境检查并提交，后续反馈用于评价或训练。该图是机制示意，不是方法优劣或学习程度的实测排名。}
\label{fig:ch06-ladder}
\end{figure}

% 图主体在 ch06-ladder-body.tex，便于单独编译与复用。
\begin{figure}[htbp]
\centering
% Mechanism map; no empirical ranking is encoded.
\begin{tikzpicture}[
 box/.style={draw,rounded corners=2pt,text width=31mm,minimum height=13mm,align=center,font=\small},
 arrow/.style={-{Stealth[length=1.6mm]},thick}
]
\node[box,fill=gray!10] (o) at (0,0) {可见观测\\候选、状态与范围};
\node[box,fill=blue!8] (p) at (4.3,0) {操作选择\\规则、拟合或训练};
\node[box,fill=orange!10] (g) at (8.6,0) {环境校验\\权限、版本与预算};
\node[box,fill=green!10] (m) at (8.6,-2.1) {状态变更\\提交并返回回执};
\node[box,fill=gray!10] (r) at (4.3,-2.1) {后续任务反馈\\可能延迟或缺失};
\draw[arrow] (o)--(p);
\draw[arrow] (p)--(g);
\draw[arrow] (g)--(m);
\draw[arrow] (m)--(r);
\draw[arrow,dashed] (r)--node[right,font=\footnotesize]{训练时可用}(p);
\node[align=left,text width=132mm,font=\footnotesize] at (4.3,-3.5)
{学习改变的是指定的决策对象；工具语义与权限仍由环境约束。虚线表示可能存在的训练反馈，不表示每次部署都会更新策略。JitRL 调整任务策略，不能据此认定其记忆操作也被学习。};
\end{tikzpicture}

\caption{记忆操作的决策边界：策略提出动作，环境检查并提交，后续反馈用于评价或训练。该图是机制示意，不是方法优劣或学习程度的实测排名。}
\label{fig:ch06-ladder}
\end{figure}

% 图主体在 ch06-ladder-body.tex，便于单独编译与复用。
\begin{figure}[htbp]
\centering
% Mechanism map; no empirical ranking is encoded.
\begin{tikzpicture}[
 box/.style={draw,rounded corners=2pt,text width=31mm,minimum height=13mm,align=center,font=\small},
 arrow/.style={-{Stealth[length=1.6mm]},thick}
]
\node[box,fill=gray!10] (o) at (0,0) {可见观测\\候选、状态与范围};
\node[box,fill=blue!8] (p) at (4.3,0) {操作选择\\规则、拟合或训练};
\node[box,fill=orange!10] (g) at (8.6,0) {环境校验\\权限、版本与预算};
\node[box,fill=green!10] (m) at (8.6,-2.1) {状态变更\\提交并返回回执};
\node[box,fill=gray!10] (r) at (4.3,-2.1) {后续任务反馈\\可能延迟或缺失};
\draw[arrow] (o)--(p);
\draw[arrow] (p)--(g);
\draw[arrow] (g)--(m);
\draw[arrow] (m)--(r);
\draw[arrow,dashed] (r)--node[right,font=\footnotesize]{训练时可用}(p);
\node[align=left,text width=132mm,font=\footnotesize] at (4.3,-3.5)
{学习改变的是指定的决策对象；工具语义与权限仍由环境约束。虚线表示可能存在的训练反馈，不表示每次部署都会更新策略。JitRL 调整任务策略，不能据此认定其记忆操作也被学习。};
\end{tikzpicture}

\caption{记忆操作的决策边界：策略提出动作，环境检查并提交，后续反馈用于评价或训练。该图是机制示意，不是方法优劣或学习程度的实测排名。}
\label{fig:ch06-ladder}
\end{figure}


\section{对照实验：E03b 能测什么，测不到什么}
\label{sec:ch06-e03b}

表~\ref{tab:ch06-who} 把决策权排开之后，下一个问题是：\emph{交出决策权这件事到底值多少}。
本书把这个问题交给 E03b。需要先说清它的级别：E03b 是\textbf{B 级受限实验}，
截至本章写作时尚未运行；本节写的是它的对照设计与受限边界，\textbf{不报告任何结果}，
也不允许从本节推出“学习化优于规则化”这类结论。

\subsection*{三臂与两条同源要求}

\begin{itemize}
  \item \texttt{rule}：写、留、取、删全部使用预先给定的规则；如果另行拟合阈值或权重，应单独登记这项学习。
  \item \texttt{self\_organize}：写入时检查近邻，再决定是否改写。两笔开销分别按检查次数和实际改写次数记录。
  \item \texttt{learned\_policy}：把记忆操作放进动作空间，由经过训练的策略选择调用。
\end{itemize}

若要把比较收窄到操作选择，三臂需要满足以下控制条件；
否则仍可比较完整方案，但不能将差异全部归于决策权：

\begin{enumerate}
  \item \textbf{检索后端固定。}三臂共用相同的匹配函数和接口；策略可选择是否调用，
        返回内容也可能因存储变化而不同。若连匹配算法也改变，需将其作为另一项干预，
        不能把整套变化仅称为“学会了取”。
  \item \textbf{任务集与留出族固定。}历史设想是复用 E03 题目作接口对照。
        这些题已经被开发与审计读取，不能在新确认性研究中重新宣称未见或未暴露。
        正式评价须另冻结来源、训练与留出边界；留出族不进入训练。
        各臂从相同存储与可见信息起步，再明确哪些后续差异属于目标干预，
        不能既允许记忆更新，又要求整个闭环始终看到相同文本。
\end{enumerate}

\subsection*{四本账，而不是一个总分}

E03b 的成本账按生命周期分成四本（表~\ref{tab:ch06-ledger}）。
分账的理由很直接：各方案的费用发生在不同位置，只报一个总数会隐藏这些差异，
看不出这笔钱到底花在哪、又从哪儿收回来。

\begin{table}[htbp]
\centering
\caption{E03b 的四本账：每一本回答一个不同的问题}
\label{tab:ch06-ledger}
\begin{tabular}{>{\raggedright\arraybackslash}p{2.2cm}>{\raggedright\arraybackslash}p{4.6cm}>{\raggedright\arraybackslash}p{6.2cm}}
\toprule
账本 & 计入什么 & 它要回答的问题 \\
\midrule
构建 & 写入一条记忆所需的调用、token 与墙钟 & 准入判决本身贵不贵 \\
维护 & 近邻重写、摘要、压缩、删除所触发的连锁调用 & 经验记忆存储变大之后，单次写入的代价怎么涨 \\
读取 & 每次检索的调用与注入上下文的长度 & 检索范围扩大，执行是否受干扰 \\
任务执行 & 完成任务本身的调用数与结果 & 前三本账花掉的钱有没有在这一栏收回来 \\
\bottomrule
\end{tabular}

\begin{minipage}{0.92\linewidth}
\vspace{2mm}
{\footnotesize 表注：四本账是本书的记账口径，不是任何一篇原文的口径；
截至本章写作时 E03b 未运行，四栏均无实测取值。}
\end{minipage}
\end{table}

不同方案可能把成本分配到不同环节，但不能在试验前规定哪一臂必然更便宜或更有效。
规则也可能调用昂贵检查，学习策略也可能在训练后只需一次轻量选择。
分项账帮助定位费用，总账回答整体代价，两者应同时报告。任何构建与训练投入都不能因发生在部署前而自动免费。

\subsection*{受限边界（必须随结果一起报）}

\begin{itemize}
  \item \textbf{动作空间是人工固定的。}AgeMem 的六个工具与 Mem-$\alpha$ 的三种写入动作
        都来自各自原文的设计\cite{agemem-2601,memalpha-2509}，
        所以 E03b 学的是“在给定工具集上何时调用”，不是“工具集本身该有什么”。
  \item \textbf{冲突消解未被独立评估。}本章反例一指出 Mem-$\alpha$ 在 Section 3.4 显式把
        Conflict Resolution 排除在评测之外\cite{memalpha-2509}；
        未设置该维度不等于删除没有奖励后果；问答正确率和操作合法性仍可能间接约束删除。E03b 应单列误删与冲突场景的检查，避免只靠一个总体分数。
  \item \textbf{训练预算与部署预算不同源。}MEM1 用 \cited{4} 张 H100 或 H200 做 PPO，
        报告的训练资源与推理资源分开列出\cite{mem1-2506}；
        E03b 的训练预算尚须在执行前冻结。若只用固定脚本或提示代替训练策略，只能算接口验证，不能标作学得策略的能力比较。
  \item \textbf{广播优势不等于逐动作因果归因。}AgeMem 的做法是把终端优势广播到
        同一轨迹的先前步骤\cite{agemem-2601}；
        它将终端结果传到先前动作，有助于训练，但不能单独证明某一步的边际贡献。分析单位应匹配任务流与持久状态的共享范围。
\end{itemize}

\section{条件变化：换任务族、换模型、换反馈器}
\label{sec:ch06-conditions}

换任务族时，应先核对各方法依赖的特征、观测与奖励。
A-MAC 的内容类型先验由对话域中的规则定义，学习过程选择它在总分中的权重，
换到代码或工具调用场景，需要检验这些定义是否仍适用；
原消融显示移除类型先验时降幅最大，未证明某种跨域失效顺序\cite{amac-2603}。

换模型时，经验文本、状态表示和工具调用方式都可能需要重新验证。
MEM1 的内部状态由被微调的策略自己生成，换基座等于换掉压缩策略本身；
A-mem 的链接与演化描述由底层语言模型生成，
原文承认组织质量受底层模型能力影响\cite{amem-2502}。
反过来，Mem-$\alpha$ 刻意把评测管线解耦成固定的 BM25 加冻结生成器，
因此换记忆策略时奖励口径不变——这是它设计上的优点，也意味着它测到的是“在这个固定口径下”的收益。

换反馈器时，问题变成奖励有没有覆盖到那个动作。
AgeMem 的奖励里含一项由语言模型充当评委的任务完成度\cite{agemem-2601}，
Mem-$\alpha$ 的内容合法性也由一个 \cited{32B} 的模型评判\cite{memalpha-2509}；
这些评分器的可靠性不能仅由任务总分确定，需要检查其输入及评价覆盖范围。

\subsection*{记忆状态非平稳：训练集不是稳定的环境}

上面三条换的都是\textbf{系统的外部条件}。还有一类变化来自系统自己，
它在固定规则与学习策略中都可能出现：\emph{经验记忆存储本身在训练过程中一直在变}。

即使更新规则固定，存储内容变化也会改变后续输入分布；学习策略又增加参数变化与数据采集之间的相互作用：
早期写入的条目决定了后来的检索结果，检索结果又决定了后来的奖励。
这会改变策略遇到的状态分布；是否构成转移规律非平稳，还需检查状态定义和外部条件，后文进一步区分。

行为监测可以检查\textbf{记忆操作的使用频率是否随训练步数变化}。
AgeMem 报了工具使用频率\cite{agemem-2601}，为检查行为提供了一项描述。
但频率表不等于随训练步数的完整曲线，也不单独区分“更会挑时机”与“操作整体减少”。
还应将操作前的状态、存储大小、调用结果和任务收益对齐，后文进一步展开。

\subsection*{任务流中的三种变化：适应、遗忘与干扰}

只看当前一组任务的进步，容易漏掉更新的影响范围。后来的任务使已经学过的任务变差，是遗忘；早先的更新使尚未学习的任务低于原基线，则是对新任务的负向干扰。两者都可能拉低整体表现，却对应不同的检查位置：前者要回访旧任务，后者要与新任务在更新前的基线比较。

Chen 等在固定 Gemini 3.6 Flash 权重的漂移任务流中比较 ReasoningBank、GEPA 与 Meta-Harness，分别检查当前组的适应、旧组的保持和未到来组的表现\cite{continual-self-evolution-2610}。全局提示或工作区改写可能覆盖早期规则；按相似度读取局部经验能收窄影响范围，但相近查询要求不同流程时，检索仍可能发生碰撞。该观察对应其 BIRD、BFCL 与 ScienceWorld 任务流及具体读取实现。

其中的工作流对照很直观：同一工具目录下，一类请求缺少必要参数，应先澄清；另一类参数齐全，应直接执行。只按主题接近来选经验，可能把前者的“先询问”带到后者，或把后者的“直接调用”带回前者\cite{continual-self-evolution-2610}。条件字段的价值就在这里：它要帮助读取者辨明动作何时适用，而不只是提高文字相似度。

据此设计持续记忆比较时，应保存中间检查点，在各自独立的状态分支上同时评估当前、已学与尚未学习的任务组。评估者可以持有未来组作参照，但题目与评分不能回流更新器。E03 冻结经验后的迁移评价没有这条持续更新过程，不能用它替代任务流中的保持与干扰测量。

\subsection*{奖励稀疏与延迟：分析单位跟着后果走}

许多记忆操作的收益晚于动作发生。写一条记忆的代价在当下，
收益要到几十步之后、甚至下一个任务里才可能兑现；
而“删”的代价可能要到被删的那条本该被用到时才暴露。
这就是长时信用分配问题在本章的具体形状。

AgeMem 的处置是把终端优势广播到同一轨迹的先前步骤\cite{agemem-2601}，
它提供了轨迹内终端反馈的训练路径，并不自动识别每次写入的贡献，也不覆盖任意跨任务延迟。
因此报告要同时明确目标和分母：任务结果按任务或任务流汇总，操作成本按操作与累计量记录。
步级平均值可以描述有效定义的步级量，只是不能替代任务级收益；当持久状态跨任务共享时，
还需考虑整条任务流的依赖，而非机械地对每个任务独立配对。

\subsection*{访问权限变化：谁能读、谁能写}

最后还要考虑一类部署条件：
记忆的读写权限不是恒定的。多智能体场景下别的智能体可能读同一份记忆，
隐私与合规约束会限制某些字段的留存，任务切换时可用范围也会变。

这一条对本章判据的冲击是直接的：\textbf{决策权交出去的同时，权限边界必须留在规则层}。
SSGM 提出的设计原则可以作为参照——它把记忆拆成不可变的情景日志 $K_{\text{ledger}}$
与可变的活动图，两层之间加读过滤与写校验两道 Gate，
写操作要先过校验才落进可变层\cite{ssgm-2603}。
策略可以在授权集合内学习写入内容和存储位置，但不能自行扩大权限或改写审计规则。
允许选择哪一层，与允许改变各层的写入权限，是两件不同的事。

\section{先认清学习对象：分数、动作与状态不是一回事}
\label{sec:ch06-object}

“记忆变得更好”可以指几种完全不同的变化。候选句子经过改写，变得更短；检索索引收录了更多条目，覆盖了过去查不到的问题；阈值在标注数据上重新拟合，减少了不合适的写入；策略经过训练，开始主动删除过时条目。这几种变化都可能影响最终回答，却不应共享一个没有宾语的“学习”。设计者首先需要说明，究竟是哪一个可保存、可冻结、可替换的对象改变了。

在 A-MAC 中，可拟合对象是式~\eqref{eq:ch06-amac} 的权重和阈值，内容类型的定义、特征计算方式和冲突处理流程仍由设计给定\cite{amac-2603}。在 A-Mem 中，新的笔记和链接可以不断出现，底层模型却不必因此进行参数更新\cite{amem-2502}。在 AgeMem 中，工具集合由框架提供，训练改变策略如何在状态下选择工具与生成内容\cite{agemem-2601}。如果只保存最终存储而没有保存阈值或策略版本，后来者就很难判断性能变化来自更多资料，还是来自不同的操作选择。

一个有用的检查方法是做一次“冻结后的互换”。设两个版本的系统分别拥有策略 $\pi_A,\pi_B$ 和记忆 $M_A,M_B$。只比较 $\pi_A(M_A)$ 与 $\pi_B(M_B)$，得到的是整体系统的差异。还需考虑在接口兼容时，让两条策略读取同一份起始记忆，或者让同一条策略读取两份记忆。这并不保证彻底分离因果效应，因为新的策略随后还会改写状态，但至少明确了差异从哪里开始。对于不同架构，强行互换可能产生无效输入，此时应承认无法完成该对照，而不是把不兼容运行的失败归因于策略能力。

同样，打分器输出一个数不意味着它经过了学习。人工写出的加权和可以是规则；在训练集上选出权重的加权和可以是拟合模型；把通用语言模型的输出直接当作效用估计，又是另一种来源。三者的公式可能一样，训练与验证的要求却不同。拟合权重需要保留选择数据和选择过程，语言模型评分需要保留提示、输入和评分结果，人工阈值需要说明适用理由。可解释的公式不会自动提供可靠的估计值。

\subsection{准入标签正确，不等于长期收益正确}

A-MAC 的准入评价把“留不留”明确写成了可以检查分类结果的问题。这个优点也带来一个必须守住的边界：标签告诉我们哪些候选按当前标注规则应当留下，却不必等于未来所有任务的价值。某条暂时不被提问的异常处理约定，可以在一次长期任务里救回关键分支；某条被标作重要的偏好，也可能在实际使用范围内从未被检索。准入指标与任务指标之间存在关系，但关系本身需要观察。

为看清两者的差别，考虑一个\textbf{本书构造的决策例子}。假定 $p$ 已经是候选应准入的校准概率，错误保留的损失为 $c_{\rm FP}$，错误丢弃的损失为 $c_{\rm FN}$，两类损失非负且和大于零，正确决策的额外损失暂记为零。那么保留与丢弃的期望损失分别为
\begin{equation}
L_{\rm keep}=(1-p)c_{\rm FP},\qquad
L_{\rm drop}=p c_{\rm FN}.
\end{equation}
在这些假定成立时，保留的条件是
\begin{equation}
p\geq \frac{c_{\rm FP}}{c_{\rm FP}+c_{\rm FN}}.
\label{eq:ch06-loss-threshold}
\end{equation}
若两类损失分别设为一和九，阈值为\calcnum{0.1}；若改为四和六，阈值变为\calcnum{0.4}。这些数是人为指定损失后的算术结果，已在随书离线脚本中核验，不是任何论文拟合出的参数。它说明：即使概率估计完全相同，任务对错误的容忍度改变，也可能合理地改变准入决策。

不能把这个推导直接套到 A-MAC 的加权分数上，宣称其阈值应该是多少。加权分数未必经过概率校准，标注误差也未必对称。式~\eqref{eq:ch06-loss-threshold} 的作用是迫使设计者说出损失假设。若关心的是长期任务收益，还要把存储费用、以后读入的上下文费用、干扰其他条目的机会成本及未来需求分布纳入。一个通用阈值常常把这些异质代价压成同一个数，因此更换任务流后需要重新检查。

更细的做法是保留分数区间上的表现。处于阈值附近的候选有多少，误判集中在哪种内容，哪些高分条目后来没有被使用，哪些低分条目被人工恢复。这些观察可以指导下一轮标注或训练。不过，“没有被使用”也不是无用的充分证据：检索器可能漏掉了它，权限可能禁止本次访问，观察窗口也可能太短。要解释一个未使用的条目，应先检查它是否有过适当的使用机会。

\subsection{把即时修正与后续复用分开测量}

StreamMemBench围绕同一证据锚安排初次任务与后续任务，分别检查证据保留、首次使用、即时反馈吸收和后续复用\cite{streammembench-2026}。其中，即时反馈吸收率FI只在收到纠正反馈的子集计分，后续复用率FUR覆盖后续任务。两列的分母不同，差值不能直接读作遗忘率。

为检验写入交互的作用，该研究从提交前的同一记忆检查点分出两支，向两支提出相同的后续请求。一支保持只读，另一支经原生接口写入完整的前段交互。研究为每个系统选取\cited{20}个证据锚，构成合计\cited{160}组配对比较。系统宏平均复用率增加\cited{13.1}个百分点\cite{streammembench-2026}。这一差值对应整段交互的联合写入，其中包含初次请求、原回答、模拟反馈与反馈后的回答。两次任务围绕同一证据锚展开，成绩由模型评审，报告值的适用范围也随这些条件限定。

\subsection{动作频率只是路径描述}

工具调用统计有助于揭示行为变化，却不能独立评定这种变化的价值。训练后删除次数减少，可能是写入质量提高，也可能是策略学会回避有惩罚的工具；检索次数增加，可能改善证据覆盖，也可能是在重复搜索已经拥有的资料。同样的频率变化需要结合操作前状态、操作结果和最终任务表现才能解释。

因此，操作统计至少需要区分提出、接收、应用和生效。策略提出了 DELETE，接口可能因版本冲突拒绝；接口接受删除后，索引可能尚未刷新；活动条目消失后，旧缓存还可能继续被读取。若只统计模型生成了多少次工具名，就会把这些不同阶段混成一次成功操作。反过来，接口为了重试而记录了两次请求，也不意味着策略作出了两次独立决策。

一个写得好的结果表应当能回答这样的问题：策略提出的更新中，有多少通过权限与版本检查；通过的更新里，有多少改变了可检索状态；这些变化发生在什么任务阶段。它不需要把所有字段塞进正文，但原始记录必须支持这些分母。否则“学会了适时操作”很容易只是一张动作柱状图的过度解释。

\section{触发门的收益要从门前开始计算}
\label{sec:ch06-gate-cost}

D-MEM 的机制值得单独拆开，因为它展示了一种经常被误读的优化：昂贵路径被少走，并不意味着被跳过的候选完全免费。其路由分数含一道效用硬阈值，过阈值后的评分由效用乘以带偏置的语义惊异程度，并设置上界。路由阈值再将输入分为跳过、只构建和深度演化\cite{dmem-2603}。低效用输入即使语义新奇，也会被硬阈值拦下。效用判断在分支选择之前就要发生。把深度演化次数作为成本代理，能够描述某条慢路径的负担，却会遗漏门前已经发生的模型调用和向量运算。

原文将该组合分数称为 RPE。这个名字提示了设计灵感，不能替代数学定义。标准时序差分误差比较即时奖励与相邻状态价值；这里的效用、惊异程度和阈值构成一条启发式路由规则。它可以有用，也可以经过实验验证，但不能仅凭名称声称它已经估计了某个策略的价值函数。阅读其他借用认知科学术语的方法时，也应先找输入、输出与更新方程。

本书再构造一组\textbf{无模型调用的记账数}。设一百个候选原本都走代价为四十单位的慢路径，总计\calcnum{4000}单位。新方案每个候选先付八单位触发门代价，其中二十个仍走慢路径，八十个走一单位的快速路径。总账为
\begin{equation}
C_{\rm gate}=100\times8+20\times40+80\times1
=\calcnum{1680}.
\end{equation}
相对旧方案节省\calcnum{58\%}，而不是由慢路径次数下降直接得出的八成。这里的单位只是教学例子的统一费用单位，不能换成 token 或金额去引用；原论文的节省比例仍应按它实际测量的接口 token 口径阅读。这个例子只证明一个算术事实：门前成本和分支成本必须使用互斥的账目相加。

实际系统还可能出现共享工作。新颖性计算得到的向量同时被后续检索使用，是否将它记为触发门成本还是构建成本，是分项口径的选择；总账不能因此多算一次，也不能两边都不算。缓存命中则涉及另一种区别：应用请求已经发生，但外部服务调用没有发生。把应用层请求次数、实际网络尝试次数和成功的模型响应次数分别记录，才有可能解释缓存和重试带来的费用变化。

触发门还会改变数据分布。被跳过的候选不再进入深度演化，后来就更少产生关于它们的收益反馈。如果训练数据主要来自被接收的候选，打分器可能对被拒绝区域越来越没有把握。可以用离线留存的候选做抽样检查，也可以在合适且获准的环境中设计探索，但不能让在线系统为收集标签任意扩大写入权限。探索机制本身的代价和风险需要属于设计的一部分。

冷启动尤其容易让“同一阈值”变成不同策略。存储为空时，新颖性高不代表内容有价值；存储稀疏时，过严的门可能导致没有任何可用经验。D-MEM在预热期将越过低阈值的输入限制为只构建笔记，不进入深度演化；这不表示所有输入都被强制存储\cite{dmem-2603}。比较触发门方法时，要同时说明从空存储起步还是从成熟存储起步、冷启动覆盖多少候选、何时切回常规规则。只比较稳定运行后的触发率，会遗漏系统形成可用记忆之前的投入。

如果要检验学习策略是否比规则门更会挑选，首先应让两个方案拥有近似的操作预算。可以在训练或开发数据上选择规则阈值，使其触发率接近学习策略，再到未参与选择的数据上比较任务结果。这样能减少“一个方案只是多做了工作”的解释。触发率匹配本身仍不意味着费用完全相同，因为被选中的候选可能长度不同、近邻不同、处理难度不同；所以结果还要回到完整资源账与任务收益。

\section{可运行的操作接口：把策略选择与状态变更分开}
\label{sec:ch06-interface}

为了让本章的决策权有一个可以动手检查的落点，随书材料提供了一个很小的离线示例。它不调用模型，不拟合权重，不计算策略优劣，也不占用 E03b 的实验编号。示例只实现 ADD、UPDATE、DELETE 和按范围读取，用来验证一个前提：当策略说“更新这条记忆”时，环境究竟承诺执行什么。

条目由内容、适用范围和版本组成。写入请求包含请求号、动作、条目标识、调用者范围、期望版本及新内容。这里的范围是测试夹具提供的已认证输入；真实系统必须从身份与权限层取得，不能相信模型自己填写的身份。操作函数依次检查动作类型、范围、存在性、版本和参数，全部通过后才改变存储。在返回回执和保存事件前，函数还会独立按条目范围隐藏前后内容，使动作名无效的越权请求同样得到脱敏回执。

回执中的 \texttt{applied=true} 表示该请求号对应的操作已生效。假设第一次更新成功后响应在途中丢失，调用者用同一请求号和相同载荷重试，仍会收到原来的成功回执，存储里却只有一次更新。统计写入次数时应按请求号去重。若同一请求号携带不同载荷，接口报错，原回执与事件记录都保留。

“版本冲突”表示请求中的期望版本与条目当前版本不符。一个常见原因是条目在读取后又被更新。收到“范围拒绝”，说明请求超出调用者的权限范围，回执与该事件都隐藏了条目的前后内容。被拒的删除应单独计数。本例在单进程内存中保存这些回执；要让进程重启后仍能识别重试，服务还需持久保存请求与状态变更的对应关系。

版本检查解决的是另一类问题。假设策略读到版本一，随后另一个合法操作将条目更新为版本二。原策略再提交“删除版本一”，系统应返回冲突，而不是删除它从未看过的新内容。策略可以重新读取再决定，但不能将冲突自动转成强制覆盖。否则看似平稳的训练回合，会在真实协作环境中变成依赖请求时序的破坏性行为。

示例还检查删除后重建的情形。若条目被删除后重建时版本又从一开始，旧请求可能误以为看到了原来的版本一，从而覆盖新对象。这个问题通常被称为 ABA 问题：标识或表面状态回到了同一个值，中间已经发生的变化却被隐藏。示例根据历史分配更大的版本号，避免这类陈旧请求被误接收；它并未实现分布式版本时钟，因此结论仅限于单个教学对象。

\begin{table}[htbp]
\centering
\caption{离线操作示例的检查范围：接口性质不等于学习效果}
\label{tab:ch06-contract}
\begin{tabular}{>{\raggedright\arraybackslash}p{4.0cm}>{\raggedright\arraybackslash}p{9.0cm}}
\toprule
检查 & 必须观察到的结果 \\
\midrule
新增与版本 & 第一次新增成功并产生版本 \\
重复请求 & 相同请求返回旧回执，不增加事件 \\
请求号被改用 & 不同载荷不能复用原请求号 \\
范围隔离 & 越权读不到内容，越权删除不改变状态 \\
陈旧版本 & 旧版本的删除被拒绝 \\
有效删除 & 活动条目消失，示例历史保留前值 \\
删除后重建 & 新版本递增，旧更新不能覆盖重建对象 \\
\bottomrule
\end{tabular}
\end{table}

七项检查在当前脚本上复算通过，结果与脚本摘要见\path{analysis/memory-operations-reproduction.json}，复算入口为\path{calculations/reproduce_memory_operations.py}。这里的“通过”只覆盖表~\ref{tab:ch06-contract} 的断言。示例没有判断一条知识是真是假，也没有测试哪种删除策略提升了后续回答率。历史保留是为了展示回放，不代表真实部署应当永久保存所有被删除内容；保留规则必须由适用的数据管理要求决定，策略本身无权把这份教学历史当作无限保存的许可。

把接口稳定下来之后，规则、拟合阈值或训练策略才能提出可比较的动作。若一条策略的 DELETE 是隐藏条目，另一条策略的 DELETE 会连同关联内容和索引一起删除，两者即使调用同一个名字，也不处在同一个动作空间。比较操作学习化，必须从这些语义约定开始；否则奖励函数学到的可能只是两个工具实现之间的差别。

HEAMS v1 的能力表注明，Mem0 原生提供按 ID 删除，评测适配器未暴露这个接口。HippoRAG 2 接受精确文本来删除条目，如何从自然语言遗忘请求定位到这段文本，仍是另一项工作\cite{heams-2611}。

谁发起删除，会改变实验的解释。论文附录 A.2 描述的协议由框架在答复前派发已识别的删除请求，更新则经普通写入完成\cite{heams-2611}。作者公开的一份实验配置却把删除工具交给模型，并在提示中要求它遇到遗忘请求时调用\cite{heams-campaign-2026}。仅凭配置无法确定表 23 中每次删除实际走了哪条路径。要判断策略是否学会何时删除，需要在同一次运行中连起三条记录：谁发起删除、接口如何回应、后续还能否读到相应内容。


\section{状态写进公式以后，策略真的看见了吗}
\label{sec:ch06-observation}

把操作写成 MDP 只是建模起点。AgeMem 的状态表达式包含上下文、长期记忆和任务描述\cite{agemem-2601}，实现时却仍需回答：策略读到的是全部长期记忆，还是一次检索返回的片段；它能否知道某条记录已经过期，能否观察一次删除引发的索引更新，能否看见下一个任务会问什么。数学上把信息放进状态，不等于工程上已经将信息交给模型。

为避免这个混淆，可以把环境状态记为 $x_t$，把策略实际收到的观测记为 $o_t$。环境可以知道全部记录及其版本，策略却只接收权限过滤后的检索结果。动作按 $\pi(a_t\mid o_t)$ 产生，环境再依照转移规则更新 $x_t$。如果当前观测不足以决定未来分布，策略就需要利用历史或内部状态；不能因为论文写了一个包含完整存储的状态元组，就默认实际策略拥有完全可观测环境。

这里有一个尤其重要的训练边界。训练阶段可能用标准答案计算任务奖励，也可能在任务对象中保存答案供评估程序调用。标准答案不应因此进入部署策略的可见上下文。审计需要分开检查环境持有字段、模型请求字段和奖励器输入字段。只读一个名为 \texttt{task} 的结构不足以证明泄漏，反过来，只说“答案用于训练奖励”也不足以证明没有泄漏；最终要检查实际发送的请求。

可见性还决定动作是否合理。策略不知道某条记录的版本，却被要求准确提交期望版本，是接口设计错误；策略不知道本次用户的范围，却需要自行猜测可读区域，是权限信息缺失。把这些失败全部计为策略能力差，会推动模型学习猜测隐藏条件。更好的做法是让环境提供完成合法决策所需的最小信息，同时保留策略不能改变的校验边界。

\subsection{存储变化不等于转移规律变化}

设固定规则是“新增候选通过校验后加入存储”。随着候选到来，存储当然不断变化，但从同一个完整状态和同一个动作到下一状态的规律可以始终相同。在这种定义下，状态访问分布发生变化，并不等于环境转移核非平稳。若直接把“存储在长大”称作“环境不再平稳”，会掩盖真正需要诊断的对象。

对策略而言，情况可能更复杂。它只能看局部检索结果，未观察到的存储变化会使相似观测对应不同后果；与此同时，训练又在改变策略参数，采样到的任务和动作分布随之变化。若外部检索服务升级、资料有效期改变或反馈器换了版本，转移或奖励规律本身也可能改变。这些来源应分别登记，才能知道要处理的是部分可观测、策略分布偏移，还是外部环境变化。

这一区分影响实验复位方式。每个任务都从同一快照开始，可以控制起始存储，却牺牲跨任务积累；让所有任务共享存储，可以考察长期更新，却使任务顺序进入处理条件。二者回答不同问题，没有一种天然更真实。报告需要写出快照来自哪里、哪些字段被重置、缓存是否清空、任务顺序是否固定，以及策略参数在整条任务流中是否继续变化。

如果比较两条自更新策略，不应让后跑的方案直接继承先跑方案留下的存储。这样会把前者的探索成果、错误与缓存一起带进后者。应从相同起点建立各自独立的状态分支，再按预先规定的任务流运行。若研究目标恰好是接管已有存储，则应将“接管哪一种来源的存储”写成显式条件，不能靠执行顺序偶然形成。

\subsection{摘要损失了什么，需要另一个观察窗口}

SUMMARY 或 FILTER 作用于工作上下文时，最容易得到的即时收益是输入变短。最难看到的代价是某条未来才有用的信息消失了。只有压缩率而没有后续任务表现的奖励，会偏爱空摘要；只有当前问答正确率，则可能允许丢弃下一问需要的条件。这些并不是某个特定算法独有的问题，而是观察窗口决定了什么能够被计入后果。

MEM1 用学习得到的内部状态压缩历史，但保留窗口与清理时点由框架限定\cite{mem1-2506}。因此，评价其状态管理应当同时看内容保留与框架约束。一个策略在固定窗口里成功处理长任务，并不能说明它已经学会任意选择存储容量；一个摘要显著变短，也不能说明被删除的信息确实无用。要把这种能力转移到新的任务长度或查询分布，需要让那些延迟使用的条件真正有机会进入评价。

\section{删除动作的信用：损失可能晚于奖励窗口}
\label{sec:ch06-delete-credit}

考虑一条存着“空栈应抛出异常”的契约记忆。当前任务只要求实现入栈，不会访问空栈。删除该条目立即减少上下文和存储，当前任务也照样成功。下一项任务才要求实现出栈，删除的代价到此才显现。若每次任务结束就结算全部奖励，并且下一项任务的失败不会回传到前一次删除，那么奖励中即使包含“任务正确率”，也可能无法约束这次错误删除。

这不意味着删除一定要有一个专属奖励项。可以延长任务流的评价窗口，可以在独立检查中测试被删信息对应的查询，可以把待删除条目暂时标为不活跃并保留恢复路径，也可以设计内容保留约束。选择取决于任务需求和可观察性。真正的问题是相关后果有没有进入可用的反馈路径，而不是奖励向量里是否出现了 DELETE 这个单词。

下面仍是\textbf{教学构造}，不属于 E03b 结果。设删除立即节省两单位费用，在观察期内有四分之一的概率需要这条记忆；如果届时缺失，将损失十单位。忽略其他影响时，删除相对保留的期望净收益为
\begin{equation}
\Delta U_{\rm delete}=2-\frac14\times10
=\calcnum{-0.5}.
\end{equation}
若奖励只覆盖当前任务，它看到的是正的两单位；若覆盖了约定的后续需求分布，期望结果却为负。差异来自结算范围，而不是计算误差。需要强调，四分之一是例子的设定，不是从少量真实任务推断出的概率；若需求分布未知，就应对不同合理取值做敏感性分析，而不是选一个让策略显得更优的数。

折扣也不能替代范围说明。将未来收益乘以折扣因子，表示目标本身如何权衡较晚发生的后果；完全不观察那些后果，则意味着估计器没有相应信息。两个方案可能都写着“长期奖励”，一个追踪后续多个任务，另一个只在当前回合终点评分，其学习信号仍可能完全不同。审稿时应沿奖励数据的实际产生路径检查，而不是停在目标函数中的求和符号。

\subsection{广播优势可以训练，但不能替代局部对照}

AgeMem 将同一轨迹的终端优势广播给先前步骤\cite{agemem-2601}。这为长程动作提供统一的结果信号，且与其整体策略训练目标相联系。问题在于，一条成功轨迹可以同时包含有帮助的检索、无必要的更新和一次侥幸无害的删除。共享终端优势并没有单独测量每个动作的净贡献。

如果想解释某次删除的边际作用，可以在允许的离线环境中，从删除之前的同一状态分出两个分支，一支执行删除，一支保留，然后继续同一组后续任务。这样的对照需要控制后续随机性、工具版本与预算，而且只能说明给定状态和后续分布下的局部影响。若删与不删导致策略走向完全不同，差异是这个动作及其后续响应的总效果，不应再拆成未经测量的“纯删除收益”。

也可以只重放已有轨迹，检查删除是否破坏了后来读取的依赖。这种做法便宜，但证据更窄：原轨迹是在旧状态下产生的，换了状态后策略未必还会提出相同查询。重放能够发现确定的接口冲突或已记录需求的缺失，却不能完整代表替代策略的闭环行为。报告应区分轨迹重放与重新执行，尤其不能把前者称为新的在线性能测试。

训练过程中还要警惕把奖励器解释当作额外真值。语言模型判断“这次删除合理”，可能只是认可了候选描述，而没有查验被删条目与未来问题的关系。若策略能看到评分提示，也可能学会写出迎合评分措辞的理由。合理的防线是保存动作前后状态与可复核结果，将理由作为诊断材料，同时用独立的任务后果或确定性约束检查关键判断。

\subsection{冲突消解不是单纯选新或选旧}

两条记忆看似矛盾时，可能存在不同的时间、主体或适用条件。一个项目约定空栈返回空值，另一个项目约定抛出异常；它们不应在共享记忆中被压成“空栈统一返回空值”。同一项目的新版本改变了约定，才可能需要将旧条目标为过时。处理冲突首先需要识别关系，之后才是 UPDATE、DELETE 或增加带条件的新条目。

因此，包含冲突场景的评价集不应只给“新事实替换旧事实”一种答案。至少应让合法更新、不同范围并存、暂时无法判断三类情形都有明确的验收规则。若环境只奖励删除旧条目，策略可能学到无条件偏爱新输入；若只奖励历史保留，又可能不断检索过时信息。评价必须告诉系统，保存原始观察与决定当前可用内容是两件可以分离的事。

SSGM 将不可变观察日志与可变活动图分开，并提出读过滤、写校验与周期性校正的设计原则\cite{ssgm-2603}。这给出了一个有用的架构参照，但该文这一部分是设计原则与研究议程，不能在本章被当作已验证的防错保证。保存日志也不自动提供正确真值：原始输入可能有误，权限可能限制恢复，观察本身也有时间范围。架构提供追查与修复的通道，具体正确性仍依赖检查。

\section{怎样使 E03b 真正比较操作选择}
\label{sec:ch06-e03b-design}

E03b 目前没有可报告的运行结果，因而最有价值的工作是消除一个常见伪对照：把手工脚本命名为 learned\_policy，再用它与另一套手工规则比较。接口可以先离线验证，学习策略必须有真实的拟合或训练过程、训练数据与冻结产物。若预算和后端不支持训练，应缩小主张为操作接口或规则门的比较，不能通过命名补齐学习过程。

首先要确定目标干预。如果研究的是同一操作集合中的调用选择，各臂就应共用工具语义、权限边界与参数校验；如果研究的是记忆的内容改写，操作选择可以固定，再比较生成内容。将内容、检索器、操作次数、预算和模型版本同时改掉，只能评价整套系统。整套系统比较有实际价值，但不应回答“决策权交给学习过程是否产生收益”这一更窄问题。

其次要确定训练与评价的时间顺序。训练数据用于拟合策略，开发数据用于选阈值、奖励权重与检查点，留出任务用于最终报告。若开发过程中多次查看留出表现再选择检查点，最后的留出分数就不再是一次独立验证。对于任务量有限的实验，应把选择次数和可见反馈登记下来，而不是假定只要文件名写着 test 就仍然未被使用。

跨任务记忆又增加了一个隔离维度：不只是题目文本不能泄漏，目标任务的答案、诊断摘要和操作后果也不能以“经验”名义提前进入起始存储。即使任务标识不同，几乎相同的模板、相同的答案片段或从目标轨迹压缩出的规则，都可能缩小真正的泛化距离。起始存储的来源清单应与任务划分一起冻结，检索结果也要保存，以便核对模型实际获得的内容。

\subsection{先做受控回放，再做独立的闭环比较}

受控回放可以给多个操作选择器同一候选、同一存储快照和同一预算状态，观察它们分别提出什么动作。它适合验证阈值边界、参数格式、权限拒绝与触发差异。若没有环境后续反馈，这种比较只证明“选择不同”，并不证明“选择更好”。本章的离线接口例子就停留在这一层。

闭环比较则让每条策略在独立状态分支上运行，当前动作改变后续可见内容，后续任务又产生新的写入机会。它才适合观察累计收益、错误积累和恢复能力。代价是各臂很快进入不同的状态，不能再简单逐步对齐。可以在任务流层面保持共同的外部输入顺序，并报告整条流的结果；不应把两条已经分叉的轨迹中同编号的第十步，当作相同状态下的两个动作。

对于持续记忆，随机化和不确定性分析也应尊重这个依赖。若多项任务共享一份存储，前一任务已经影响后一任务，就不能把所有任务当作彼此独立的样本来人为缩窄区间。可以从独立初始化的多条任务流得到重复，或在设计允许时按任务族及流分组分析。若只有一条流，就如实说明无法可靠估计流间波动，不用成千上万条工具事件冒充独立重复。

\subsection{先冻结停止规则，再观察差异}

失败何时停止，是策略比较的一部分。一个方案遇到拒绝后无限重试，另一个方案立即退回规则路径，两者的实际预算就不同。需要预先规定每个任务的模型调用、工具调用和修复上限，并说明到限后如何记失败、如何处理尚未完成的更新。更新如果跨越多个步骤，停止时应留下可识别的未提交状态，不能既算作新知识又不计其未完成的费用。

同样要冻结缺失值处理。模型调用失败、结果文件缺失与合法的零次操作不是一回事。账本写零可能表示确实没有发生费用，也可能只是恢复运行后忘记读回旧账。第5章 E03 的补充审计已经展示了这一风险。E03b 应保存单个事件的资源使用，再由事件合计出阶段总账；恢复运行时先核对已完成事件，避免重复计费或把既有成本清零。

最后，学习成本和部署成本应在明确的复用范围内合并。训练一次在很多任务上复用，与为每个新任务重新训练，是两种生命周期。若训练耗费为 $C_{\rm train}$，在 $n$ 个约定任务上报告平均总成本，可以写成 $C_{\rm train}/n+\overline C_{\rm deploy}$，但必须说明 $n$ 是实际完成的任务数，还是计划中的复用假设。不能用尚不存在的大量未来任务把当下训练成本摊成接近零。

\section{从实验策略走向可维护系统}
\label{sec:ch06-maintainable}

一条经过训练的记忆策略仍然需要外部控制。谁能调用工具、哪些存储允许修改、版本冲突怎样处理，这些边界属于环境契约。策略可以在合法动作中选择更合适的一项，不能通过生成新的工具参数修改自己的权限。否则性能优化和规则变更发生在同一个闭环里，结果很难解释，也难以回滚。

策略版本与存储版本应该分开保存。回退策略参数不会自动撤销它已经写入的内容；回退存储也不会自动恢复此前的检索缓存或候选评分。一次真实回退至少要说明恢复到哪份策略、哪份活动存储、哪种索引与哪个权限配置。若只将模型切回旧检查点，却继续读取新策略产生的记忆，运行的是一个混合版本，应当按混合版本重新检查。

保守回退可以先限定操作范围。例如，当新策略在测试中出现异常删除时，暂停其删除提案的自动应用，保留只读检索和必要的新增；当评分器失效时，可以使用经过验证的固定规则处理明确场景，将模糊候选留待后续处理。这样的设计并不保证性能不下降，它提供的是可解释的降级行为。任何回退都应记录发生条件及期间累计的待处理内容，以免恢复后把所有积压候选一次性无条件写入。

监测也应围绕具体失败模式。权限拒绝率上升可能意味着调用上下文变化，版本冲突增多可能来自并发更新，摘要长度突然下降可能来自奖励偏好，删除后恢复率上升则可能意味着保留价值被低估。这些信号值得检查，但不应自动当作训练失效的证明。先沿请求、回执和状态变化定位原因，再决定是修接口、调规则还是重新训练。

本章给出的选择顺序由此变得清楚：先明确操作语义和合法边界，再检查策略可见的信息是否足够；用离线回放排除接口错误，在冻结预算与起始状态后比较闭环结果；最后才讨论哪些决策值得持续交给学习过程。读者不必为每个记忆系统训练一个新模型，但应能指出自己正在依赖哪一种学习、哪一种规则，以及哪一类后果还没有进入评价。


\subsection{先核对汇总，再判断比较范围}
\label{sec:ch06-summary-check}

方法综述需要先确认汇总值怎样得到。Mem-$\alpha$ 的作者在官方仓库说明，Avg.\ 应按数据集等权平均，并将验证集 Table 1 的训练版本均值由\cited{0.642}更正为\cited{0.619}；更正后的测试集 Table 2 仍为\cited{0.592}\cite{memalpha-erratum-2025}。这是对已报告表格的更正，没有报告一次新的训练运行。

\begin{table}[htbp]
\centering
\caption{Mem-$\alpha$ 的汇总值：原版、作者更正与显示列复算}
\label{tab:ch06-memalpha-averages}
\begin{tabular}{lccc}
\toprule
对象与位置 & v1 均值 & 作者更正 & 按显示列复算 \\
\midrule
验证集训练版（Table 1） & \cited{0.642} & \cited{0.619} & \calcnum{0.6190} \\
测试集训练版（Table 2） & \cited{0.592} & \cited{0.592} & \calcnum{0.5921} \\
验证集基座（Table 3） & \cited{0.389} & 未单列 & \calcnum{0.3914} \\
\bottomrule
\end{tabular}

\vspace{2mm}
\begin{footnotesize}
原版见\cite{memalpha-2509}；作者更正只重列 Table 1、2，见\cite{memalpha-erratum-2025}。复算对已显示的任务分数等权取平均，使用的是舍入后的列值，不替代未舍入结果。Table 3 的基座均值尚无单列更正，本书不据原显示均值补报精确训练增益。
\end{footnotesize}
\end{table}

验证集与 MemoryAgentBench 测试集由不同任务组成，指标也按任务使用 F1 或准确率；不能把前者的基座均值从后者的训练版本均值中相减，声称得到同一测试集上的训练收益。表~\ref{tab:ch06-memalpha-averages} 将汇总检查与比较范围分开：先问数字能否由对应列恢复，再问两列是否来自同一组实验条件。

这些差异也会改变消融实验的读法。比较原论文表4中内容奖励系数同为 $\gamma=0.1$ 的两行，压缩奖励系数 $\beta$ 分别为零和 $0.05$。均值栏依次给出 \cited{0.630} 和 \cited{0.642}。若按七个数据集等权复算，两行显示成绩的均值依次为 \calcnum{0.61957} 和 \calcnum{0.61900}。它们相差不到千分之一；各列仅显示到小数点后三位，未舍入成绩的排序仍待核。对于 BookSum，两行的记忆长度则从 \cited{4.5} 千词元降至 \cited{2.2} 千词元。读者可以据此比较这项任务上的长度变化，不能从原均值栏认定哪种系数使总体成绩最高\cite{memalpha-2509}。

读取平均值时还要检查各列是否共享量纲。AgeMem 报告的跨数据集平均表现汇集了成功率、进度率与评委分数\cite{agemem-2601}。这种汇总可以忠实复述为论文所用的平均表现，但不宜改名为所有任务的统一“成功概率”。同理，MemGen 摘要中的最大增益只能按其对应基准和比较方法理解，不能与另一篇论文的跨任务均值直接排序。

\subsection{默认值要沿启动配置追到计算位置}
\label{sec:ch06-config-identity}

同一个系数可能同时出现在通用默认项和具体训练参数里。Mem-$\alpha$ v1 的第7页实现说明写主实验 $\gamma=1$，第8页消融说明写默认 $\gamma=0.1$，两处都配合 $\beta=0.05$\cite{memalpha-2509}。在第9页表4，标作 $(\beta,\gamma)=(0.05,0.1)$ 的一组数逐列重现了表1七个验证数据集的成绩和记忆长度。第21页表9用同样的系数标注，逐列重现了表2中4B模型在九个测试数据集上的这两项读数\cite{memalpha-2509}。两张消融表对默认行的标注彼此一致，第7页对主实验的说明仍有分歧。固定公开提交的主训练脚本则明确设置压缩系数\cited{0.05}、内容系数\cited{0.1}；训练入口将它们传给奖励管理器，后者实际乘到对应奖励项上\cite{memalpha-code-2025}。按这份脚本的显式参数，通用配置中的\cited{1.0}回退值不会自动成为此次启动的内容系数。

KL 约束还要沿两条路径检查。原文描述去掉 GRPO 的 KL 项以鼓励探索\cite{memalpha-2509}；这份公开主脚本沿用默认配置关闭奖励中的 KL 项，并显式开启系数为\cited{0.001}的策略 KL 损失，策略更新代码会将它加进损失函数\cite{memalpha-code-2025}。只检查奖励侧开关，会漏掉策略损失侧的约束。

这些代码能够说明固定脚本怎样传递参数，尚未将其绑定到原表使用的训练检查点。复现报告结果需要同时保存训练命令、合并后的有效配置和实际检查点；解释原文目标与解释公开实现时，也应分别标明来源。这样才能区分方法描述、当前可读实现和已经完成的历史运行。


% ---------- 6. 选择原则 ----------
\begin{principlebox}
\textbf{先问决策权，再问算法。}
把“写、留、取、删”四个决策逐个列出，标出每一个的决策者是人、是规则还是学习过程；
只对那些确实值得交出的决策付训练成本。

\textbf{交出决策权必须同时交出账。}
每一种可能有害的操作都应有适当检查或反馈，不必机械地为每个动作新增一个奖励项。
AgeMem 报了工具使用频率\cite{agemem-2601}，
这提供了行为描述；复核成本与后果还需要事件记录、分母和测量条件，不能由频率表直接推定可复现。

\textbf{规则门是强化学习的对照，不是它的下位替代。}
D-MEM在所报告配置下取得约\cited{80\%}的token降幅\cite{dmem-2603}，
对具有相容输入和操作接口的学习策略，可以加入触发率相近的规则门对照，
报告相同预算约束下的收益。D-MEM 的整套流程并非适配所有架构的通用基线。

\textbf{潜记忆要单独定审计规则。}
MemGen 的记忆是定长潜向量\cite{memgen-2509}，
“这条记忆写了什么”无法读出；
它仍可通过外部行为与资源检查，但不能提供自然语言条目式的内容审计；引用收益时应区分这两种可验证性。
\end{principlebox}

% ---------- 7. 反例 ----------
\begin{pitfallbox}
\textbf{反例一：动作空间里有“删除”，评测维度里却没有冲突消解。}
Mem-$\alpha$ 的动作集包含 \texttt{memory delete}，
而原文在 Section 3.4 显式把 Conflict Resolution 排除在评测之外，
理由是缺少可用的真实基准\cite{memalpha-2509}。
于是“删掉一条与未来问题冲突的旧记忆”这个动作，
在训练信号这一侧既可能记正收益，也可能零成本——
下游正确率与语义合法性可能给出间接惩罚，但原评测没有单列冲突消解，不能据此声称这项能力已得到覆盖。
这正是第 14 章要展开的治理问题。
SSGM 的设计建议是把记忆拆成两层：不可变的情景日志 $K_{\text{ledger}}$ 与可变的活动图，
在两层之间加读过滤与写校验两道 Gate，写操作要先过校验才落进可变层\cite{ssgm-2603}。
这条设计的代价是“删除”不再是一个可由策略自由调用的原语——它先要回答“删的是哪一层”。

\textbf{反例二：把“用了 RL”当成“记忆操作被学习了”。}
JitRL 通过优势估计调制 logit，写入与检索采用固定规则；本书未在原文定位到显式遗忘机制，不能把它说成已经实现的固定遗忘规则\cite{jitrl-2601}。
把任务策略的调制直接归为记忆操作学习，会让本章的判据失去区分作用。

\textbf{反例三：把“学出来的权重”当成“通用的权重”。}
A-MAC 学出的权重带来的全部增益，在数值上等于移除内容类型先验的消融差值\cite{amac-2603}。
两条差值相等不等于因果贡献相等；换域后是否需要调整先验与权重，须另行检验。
\end{pitfallbox}

\section*{习题}
\begin{enumerate}[label=6-\arabic*.,ref=6-\arabic*]
  \item\label{ex:ch06-1} 选一个你正在用的记忆系统，按表~\ref{tab:ch06-who} 填出四个决策的决策者。
        如果四栏全是“规则”，估算一下把它换成学习策略后，新增的可观测账目有哪几项。
  \item\label{ex:ch06-2} 式~\eqref{eq:ch06-amac} 里 $U$ 用了语言模型而其余四项用规则。
        分别把 $U$ 换成规则、把 $T$ 换成语言模型打分，比较语义判断、可复核性和费用可能怎样变化。还需哪些观察，才能判断准入效果是否改善？
  \item\label{ex:ch06-3} 本章反例一指出冲突消解未被独立评估。
        请为 Mem-$\alpha$ 的奖励设计一项能覆盖它的分数，并说明这项分数本身可能被怎样利用。
  \item\label{ex:ch06-4} D-MEM 报告了演化触发率这道 Gate 的取法。
        若一个用强化学习学到的策略声称学会了“何时更新”，你会要求它报出哪三个数来与这道 Gate 的基线对比？
\end{enumerate}

\section*{常见误解}
\begin{itemize}
  \item \textbf{“记忆操作学习化了就不用管准入规则了”。}
        权限、接口与外部校验仍需要明确。A-MAC 的特征定义由设计给定、权重由数据选择；
        Mem-$\alpha$ 的奖励设计约束策略训练，两者不能都简写成“规则变成奖励”。
  \item \textbf{“召回高说明记得全，是好事”。}
        召回 \cited{1.000} 配精确率 \cited{0.371}\cite{amac-2603}，
        它表示相对准入标注存在较多假阳性；是否损害后续使用，还需下游检查。
  \item \textbf{“无梯度方法不算学习”。}
        JitRL 不更新任何权重，但它把经验变成了策略偏置\cite{jitrl-2601}；
        本章的分界是决策权归属，不是有没有梯度。
\end{itemize}

\section*{进一步阅读}
\begin{itemize}
  \item \cite{amac-2603}：把准入写成式~\eqref{eq:ch06-amac} 的出处，
        Table 1、Table 2 与 Table 5 分别提供总体、消融与分域数字；本文已指出跨表口径待核。
  \item \cite{amac-code-2026}：固定提交的参考代码；重点对照新颖性截断、空库约定、对话级划分与准入接口，不能将代码发布等同于原表复现。
  \item \cite{continual-self-evolution-2610}：在漂移任务流上分别观察适应、保持与干扰，说明冻结模型权重以后仍需追踪更新的作用范围；结果限于其预印本中的配置。
  \item \cite{amem-2502}：自主组织的机制与 token 账，
        它的消融（Table 3）只去掉了模块，没有隔离“链接结构”与“记忆内容”，读它时注意这个缺口。
  \item \cite{agemem-2601}：六工具动作空间与三阶段训练，
        它的工具使用统计（Table 3）是“交出决策权必须交出账”的具体样子。
  \item \cite{memalpha-2509,memalpha-erratum-2025,memalpha-code-2025}：把记忆架构与 RL 解耦的设计；表格均值还应结合作者更正，训练配置则沿固定脚本核对，
        Section 3.4 关于 Conflict Resolution 的说明是本章反例一的来源。
  \item \cite{mem1-2506}：用内部状态整合记忆与推理，
        Table 4 里 RL 与 SFT 的对照说明“崩溃”在长程任务上是什么形状。
  \item \cite{dmem-2603}：规则门这一支，
        它提出的带噪基准是检验“何时演化”类主张的可复用工具。
\end{itemize}
